
\chapter[System Design]{System Design}{System Design}\label{sd}

\noindent
\rule{0.49\textwidth}{0.75pt} $_{\bigcirc}$ \rule{0.49\textwidth}{0.75pt}\\

\section{Design Diagrams}

\subsection{Data Flow Diagram}
 \begin{figure}[h!]
     \centering
     \includegraphics[width=1\textwidth]{Chapter1/DFD.png} % Adjust the width as needed
     \caption{Data Flow Diagram}
     \label{f1} 
\end{figure}
The Data Flow Diagram for DAVS represents the end-to-end flow of information in a secure federated learning round. The primary entities are the hospital clients (local trainers), the coordinator (aggregation and orchestration), the selected validator committee (consensus), and the permissioned blockchain ledger (audit).

The flow begins when each hospital performs local training on its private dataset and computes a model update. Instead of transmitting full gradients, the client applies Random Projection-based gradient sketching to generate a compact sketch vector. These sketches are sent to the coordinator, reducing communication overhead while preserving similarity structure needed for trust evaluation.

Next, the coordinator executes Data-Aware Verifier Selection (DAVS) over the collected sketches. DAVS computes round-local representativeness scores using cosine similarity and magnitude consistency, ranks clients, and selects the top-\(k\) nodes as a validator committee. Low-scoring outliers are treated as suspicious and are excluded from the validation set.

The committee then performs PBFT-inspired validation on the aggregated update. Validators run deterministic quality checks (e.g., numerical validity and bounded change) and vote to accept or reject the update. If quorum is reached, the global model is updated; otherwise, the previous model is retained. Finally, the coordinator writes per-round metadata (model hash, committee IDs, DAVS scores, PBFT votes, and evaluation metrics) to the permissioned blockchain, providing immutable provenance and auditability of the full training lifecycle.





\subsection{Use Case Diagram}


\begin{figure}[h!]
    \centering
    \includegraphics[width=1\textwidth]{Chapter1/use_case_new.jpeg} % Adjust the width as needed
    \caption{Use Case Diagram }
    \label{f1} % Optional: for referencing the figure in your document
\end{figure}
The Use Case Diagram for DAVS describes interactions among the major actors involved in secure federated medical model training. The primary actors are Hospital Clients (local trainers), the Coordinator/Aggregator (round orchestration and aggregation), Validator Committee Nodes (selected verifiers), and the Permissioned Blockchain Ledger (audit).

Hospital clients participate by receiving the current global model, training locally, producing gradient sketches, and submitting those sketches to the coordinator. The coordinator performs aggregation and runs DAVS to select a validator committee based on current-round sketch statistics. The validator committee validates the proposed global update using PBFT-inspired voting and deterministic checks, producing an accept/reject decision.

Upon acceptance, the coordinator updates the global model and broadcasts it for the next round; upon rejection, the previous global model is retained. In both cases, the ledger stores the round metadata (model hash, committee IDs, DAVS scores, PBFT votes, and evaluation metrics). This use-case flow emphasizes privacy (no raw data sharing), robustness (malicious update filtering and consensus validation), and auditability (immutable logging).


\subsection{Data Structure Design and Smart Contract}
The DAVS audit layer replaces the need for trusting ephemeral logs by using a permissioned blockchain-style ledger to store immutable round provenance. Rather than storing raw training data or full model artifacts, the data structure focuses on compact metadata sufficient for verification and audit.

Each round records: the round number, the global model hash, selected committee identifiers, DAVS scores (or summaries), PBFT vote logs, and evaluation metrics (e.g., accuracy/loss). This structure supports post-hoc forensic analysis, enables verification that a given global model corresponds to an accepted consensus outcome, and provides transparency about which nodes participated in validation.

While the implementation uses consensus and ledger concepts rather than user-facing smart contract access control, the design principles are similar: critical state transitions (accept/reject) are logged immutably, and the stored hashes ensure integrity without requiring storage of sensitive artifacts on-chain.


\rule{0.49\textwidth}{0.75pt} $_{\bigcirc}$ \rule{0.49\textwidth}{0.75pt}\\
